\section{工具设计的最佳实践}

文章的附录二是一个极具实操价值的章节，讨论了如何为 Agent 设计工具接口——这是决定 Agent 效果的关键因素。

\subsection{Agent-Computer Interface (ACI)}

Anthropic 提出了一个重要概念：\textbf{ACI}（Agent-Computer Interface），与 HCI（Human-Computer Interface）相对应。核心观点是：设计 Agent 工具接口时，应该投入\textbf{与设计人机界面同等}的精力。

\begin{importantbox}{ACI 设计原则}
\begin{enumerate}
    \item \textbf{换位思考}：站在模型的角度，仅凭工具描述和参数，能否直观地知道如何使用这个工具？如果你自己需要仔细思考，模型也一样。
    \item \textbf{完善文档}：好的工具定义应包含使用示例、边界情况、输入格式要求、与其他工具的区别。
    \item \textbf{迭代测试}：在 Workbench 中用大量示例输入测试模型如何使用工具，根据观察到的错误迭代优化。
    \item \textbf{防错设计（Poka-yoke）}：修改参数设计使模型更难犯错。
\end{enumerate}
\end{importantbox}

\subsection{工具格式的选择}

同一个操作往往有多种描述方式。例如文件编辑可以用 diff 格式或重写整个文件；结构化输出可以用 Markdown 或 JSON。虽然在软件工程中这些差异是表面的（可以无损转换），但对 LLM 来说，不同格式的\textbf{难度差异巨大}。

\begin{table}[H]
\centering
\begin{tabular}{>{\bfseries}p{3cm} p{4.5cm} p{4.5cm}}
\toprule
设计维度 & 推荐做法 & 需要避免的做法 \\
\midrule
思考空间 & 给模型足够的 token 来「思考」再写入关键输出 & 要求模型在写出内容前就确定元信息（如 diff 的行数） \\
格式自然度 & 使用模型在互联网文本中常见的格式 & 使用罕见或高度结构化的自定义格式 \\
格式开销 & 最小化格式负担 & 要求精确计数千行代码或对代码做 JSON 转义 \\
路径规范 & 使用绝对路径 & 使用相对路径（容易因目录切换出错） \\
\bottomrule
\end{tabular}
\caption{工具格式设计的推荐做法与需要避免的做法}
\end{table}

\begin{knowledgebox}{来自 SWE-bench 的实践经验}
Anthropic 在构建 SWE-bench Agent 时，在\textbf{工具优化上花费的时间远超整体 prompt 优化}。一个典型案例：模型在使用相对路径的工具时频繁出错（因为 Agent 会切换工作目录），将工具改为强制要求绝对路径后，模型使用完美无误。这说明好的 ACI 设计可以从根本上消除某类错误。
\end{knowledgebox}

\subsection{本章小结}
工具设计是 Agent 成功的关键因素，甚至比 prompt 优化更重要。ACI 的设计应遵循换位思考、完善文档、迭代测试、防错设计四大原则。工具格式应选择模型容易生成的自然格式，避免不必要的格式开销。


